% element: 10-s075:e03
% element: 10-s075:e04
\BeleskeID{10-s075}
Standardna oznaka počinje numeričkim prefiksom 74 ili 54. Slovni deo zatim određuje familiju. Oznaka F je Fast, koju je koristio Fairchild, i u ovom poređenju suštinski odgovara ALS familiji.

Naredne dve do pet cifara određuju tip kola. U njima nema opšte „suvisle logike“ za dekodovanje, pa za konkretnu funkciju treba pogledati katalog proizvođača.

Na crtežu su različiti pogledi DIP14 kućišta, razmak pinova i dimenzije u milimetrima sa vrednostima u inčima u zagradama. Razmak pinova je $2.54\,\mathrm{mm}$ odnosno $0.100$ inča, T. P. (See Note A). Kućište potpada pod JEDEC TO-116 i EIA MO-001AA dimenzije. Oznake broja mesta i ograničenja MIN/MAX čitaju se uz odgovarajuće dimenzije, a ne kao oznake električnog priključka.

\subsection*{Dimenzije kućišta}
Vrednosti u zagradama su inči. Čeoni pogled daje rastojanje osa pinova $7.62\pm0.25$ mm $(0.300\pm0.010)$, širinu $6.35\pm0.25$ mm $(0.250\pm0.010)$, širinu udubljenja $2.0$ mm $(0.080)$ NOM i njegovu dubinu $0.25$ mm $(0.010)$ NOM. Ugao pinova je $90^\circ$–$105^\circ$ na svih 14 mesta; debljina je $0.25$–$0.36$ mm $(0.010$–$0.014)$, See Notes B and C.

Dužina tela je $18.0$–$19.8$ mm $(0.710$–$0.780)$. Zarez ima nominalni poluprečnik $2.4$ mm $(0.093)$ i nominalnu širinu $2.8$ mm $(0.110)$. Bočni pogled daje visinu do ravni naleganja najviše $5.08$ mm $(0.200)$, širinu ramena pina najviše $1.78$ mm $(0.070)$ na svih 14 mesta, rastojanje donje površine tela od ravni najmanje $0.51$ mm $(0.020)$ i dužinu pina ispod ravni najmanje $3.17$ mm $(0.125)$.

Rastojanje kraja tela od ose krajnjeg pina je $2.03\pm0.51$ mm $(0.080\pm0.020)$ na četiri mesta. Prelazni deo pina ima najmanje $0.84$ mm $(0.033)$, a širina donjeg dela je $0.381$–$0.533$ mm $(0.015$–$0.021)$ na 14 mesta, See Notes B and C. Sve dimenzije odnose se na odgovarajući pogled, ne na električno ponašanje kola.
